\subsection{章末確認問題}

\begin{enumerate}
  \item ソフトウェア工学経済は、なぜ単なる「費用削減」の話ではないのか。
  \item 提案、代替案、何もしない代替案の違いを説明せよ。
  \item キャッシュフローを時間軸で表す理由を説明せよ。
  \item お金の時間価値を無視すると、どのような判断ミスが起こるか。
  \item 許容最低収益率が機会費用を表す理由を説明せよ。
  \item 営利組織と非営利組織では、意思決定技法がなぜ異なるのか。
  \item 損益分岐分析がクラウドサービス選択に役立つ例を説明せよ。
  \item 補償型技法と非補償型技法の違いを説明せよ。
  \item 無形資産を考慮しないソフトウェア提案には、どのようなリスクがあるか。
  \item 専門家判断、類推、分解、パラメトリック見積りの違いを説明せよ。
  \item 総所有費用が、初期開発費だけの比較より重要になる理由を説明せよ。
  \item 効率、有効性、生産性の違いを、自分の言葉で説明せよ。
\end{enumerate}

\begin{pointbox}{重要}
解答の方向性 1 は価値、リスク、無形資産、組織目標を含めて説明する。2 は、単一の行動案、相互排他的な選択肢、候補群を選ばない選択肢として区別する。5 は、ある案に資源を投じると別の投資機会を失う点を述べる。8 は、ある基準の不足を他基準で補えるかどうかで区別する。12 は、効率が資源消費、有効性が目的達成、生産性が価値と資源の比であることを述べる。
\end{pointbox}
